% Einheit 1 von 5 – Matrizen als Rechenobjekte (bank/matrizen-und-uebergangsprozesse/e1.jsonl, zusammenbau v0.1)
%% TODO Zeitmarke Einheit 1: Marken-Zeile nicht lesbar
%% TODO Zweigzeile Teil 1: was hier gelernt wird (Satz in Schülersprache aus dem Einheitstitel) – Einheit 1
\einheitenkopf[e1]{Einheit 1 von 5 · Matrizen als Rechenobjekte}
\zweigzeile{keine Prüfungsaufgabe}
\setcounter{aufgabe}{8}

%% TODO Titel: Ich-kann-Satz fehlt in der Bank (Platzhalter: Kettenname) – Nr. 9
%% TODO Anweisung: ein Satz über der ersten Teilaufgabe fehlt in der Bank – Nr. 9
\begin{aufgabe}{Matrizen als Rechenobjekte}
%% TODO \rechenplatz{n}: Schrittzahl des Grundfalls fehlt in der Bank (nur bei fester Schreibform, 2.3 b)
\begin{teile}
\teil Rechnen oder deuten? Nur ankreuzen, nicht ausführen: „Ermittle das Produkt der Matrix $((2 | 5), (1 | 4))$ mit dem Vektor $(3 | 2)$.“ \\ \kreuz{rechnen} \\ \kreuz{deuten}
\teil Berechne $A \cdot v$ für $A = ((2 | 1), (3 | 4))$ und $v = (1 | 2)$. ( \leerfeld | \leerfeld )
\teil A hat 2 Zeilen und 3 Spalten, B hat 3 Zeilen und 4 Spalten. Ist $A \cdot B$ bildbar, und welches Format hat es? Ist $B \cdot A$ bildbar?
\teil Berechne $M^2$ für $M = ((1 | 2), (2 | 1))$. (( \leerfeld | \leerfeld ), ( \leerfeld | \leerfeld ))
\teil Bestimme die Inverse von $A = ((2 | 1), (1 | 1))$ über den Ansatz $A \cdot B = E$ mit $B = ((a | b), (c | d))$. (( \leerfeld | \leerfeld ), ( \leerfeld | \leerfeld ))
\teil Beschreibe die Wirkung von $P = ((0 | 1), (1 | 0))$ auf einen Vektor $(a | b)$.
\end{teile}
\begin{gleichungsraster}[2]
\gl{\text{$M = ((a | 1), (2 | b))$ und $M \cdot (1 | 2) = (5 | 10)$. Bestimme a und b.}} &
\end{gleichungsraster}
\begin{teile}
\teil $A = ((0 | 0), (1 | 0))$ und $X = ((a | b), (c | d))$ mit reellen Zahlen a, b, c, d. Ermittle alle a, b, c, d, für die $(X + A)^2 = X^2 + 2 \cdot A \cdot X + A^2$ gilt \hfill (Abitur 2025 LK).
\end{teile}
\end{aufgabe}

%% TODO Titel: Ich-kann-Satz fehlt in der Bank (Platzhalter: Kettenname) – Nr. 10
%% TODO Anweisung: ein Satz über der ersten Teilaufgabe fehlt in der Bank – Nr. 10
\begin{aufgabe}{Matrizen als Rechenobjekte – Fehler finden, Begründen}
\begin{teile}
\teil Lena berechnet das Quadrat von $M = ((2 | 1), (0 | 3))$ und schreibt $M^2 = ((4 | 1), (0 | 9))$. Finde den Fehler und rechne richtig.
\teil Begründe ohne Rechnung, warum $A \cdot B$ und $B \cdot A$ bei Matrizen verschieden sein können.
\end{teile}
\end{aufgabe}
